
The field of quantum matter encompasses a diverse array of materials and systems whose properties are governed by the principles of quantum mechanics. In these systems, the behavior of particles is described by quantum theory, leading to phenomena such as superposition, entanglement, and coherence. Quantum matter includes exotic states of matter such as quantum spin liquids, superconductors, and topological insulators.

The study of matter and the phase transitions between different states of matter.

%On quantum matter, a general rule to determine a phase of matter is to use an order parameter. This parameter is the expectation value of a local observable in the given quantum state. 

%difference to long-ranged order parameter

The field of quantum spin liquids (QSLs) has seen significant development over the years, characterized by the exploration of highly entangled ground states in spin systems. Quantum spin liquids represent 'quantum disordered' phases where strong zero-point fluctuations prevent conventional magnetic long-range order. Moreover, these states exhibit massive many-body entanglement, distinguishing them as distinct phases of matter with unique physical properties, such as non-local excitations and topological features.

The term QSL was introduced by Anderson in 1973 in his seminal paper describing a 'resonating valence bond' (RVB) state as the ground state of an antiferromagnet. This RVB state, envisioned at the wavefunction level, paired spins into singlets and formed a quantum superposition of different pairings. Although originally proposed for the spin-1/2 Heisenberg antiferromagnet on the triangular lattice, this concept highlighted the role of frustration in stabilizing QSL states.

The interest in QSLs experienced a resurgence in the late 1980s, driven by the discovery of high-temperature superconductivity. Anderson and others proposed specific forms of RVB states, suggesting that superconductivity might emerge from doped QSLs. Additionally, the notion of 'pseudo-fermions', or 'spinons', as possible quasiparticles of QSLs was introduced, marking the beginning of modern 'parton' methods.

Parallel developments in the late 1980s connected QSLs to the fractional quantum Hall effect (FQHE) and topology. Kalmeyer and Laughlin proposed a 'chiral spin liquid' state, while Kivelson, Rokhsar, and Sethna suggested the role of topology in RVB states. Furthermore, the quantum dimer model was introduced to elucidate the topological structure of QSLs.

%Despite these advancements, skepticism persisted within the scientific community regarding the feasibility of the RVB paradigm in physical systems. However, 


According to  Landau's theory, patterns in phases of matter actually emerge from symmetry breaking. Without symmetry breaking, as is the case in liquids, there would be no patterns and the state would be featureless. But, in the late 1980s, it became clear that even disordered liquids can have features. In an effort to describe the exotic properties of superconductors at high critical temperatures, a two-dimensional (2D) disordered spin liquid - i.e., a chiral spin liquid—was discovered (2,3). A chiral spin liquid is characterized by the absence of any spin order and perfect edge heat conduction (because all edge excitations move in the same direction). It was quickly realized that there can be many different chiral spin liquids with exactly the same symmetry (4) but different numbers of heat-conducting edge modes (5). Therefore, symmetry alone is not enough to characterize and distinguish different chiral spin liquids. This means that chiral spin liquids must exhibit another kind of order beyond the usual symmetry description. The proposed new kind of order was named ``topological order''(4) [the term topological was motivated by the low-energy effective theory of the chiral spin liquids (3), which is a topological quantum field theory (6)]

For gapped systems: a topologically ordered state is a gapped ground state that cannot be continuously deformed into a product state without closing the energy gap (i.e., without a phase transition). Two topologically ordered states belong to the same phase (i.e., have the same topological order) if they can deform into each other smoothly without a phase transition.

the ground-state degeneracy (4,15) on a torus (and other spaces with nontrivial topology), (ii) the non-Abelian geometric  phase  (16)  of  those  degenerate ground states, and (iii) the gapless edge modes


around 2000, QSLs experienced a renewed interest, fueled by theoretical connections between spinons, fractional quantum numbers, and topological order.
Key developments during this period included the realization of stable QSL phases in lattice gauge theories and the identification of topological phases in certain QSL states. Kitaev's toric code model and honeycomb model provided simple yet exact solutions for QSL phases, dispelling doubts about their existence.

As theoretical developments flourished, experimental efforts intensified in the search for QSLs, leading to numerous discoveries. The synthesis of theoretical and experimental findings has advanced our understanding of QSLs and their potential technological applications.

Despite decades of research, the precise definition of a QSL remains elusive. However, the essential feature of a QSL is believed to be the presence of an anomalously high degree of entanglement, distinguishing it from conventional ordered phases.

Overall, the field of quantum spin liquids continues to evolve, driven by the quest to unravel the fundamental properties of highly entangled quantum matter and its implications for both theoretical physics and experimental investigations.



At the heart of quantum spin liquids lies a profound degree of entanglement, which manifests through three key features: topological order, long-range entanglement, and fractionalization. Topological order is characterized by a ground state degeneracy that eludes distinction through local measurements, thus rendering it resilient against errors. Long-range entanglement, on the other hand, is indicative of a non-zero topological entanglement entropy, offering insights into the global properties of the system. Furthermore, excitations within quantum spin liquids carry fractional quantum numbers, with two-dimensional systems often exhibiting non-Abelian anyons, paving the way for exotic quantum phenomena \cite{Kitaev2006}.

An important model in the study of quantum spin liquids is the Kitaev model, an exactly solvable spin-1/2 model on a honeycomb lattice, which serves as a prototype for understanding the fundamental properties of these intriguing phases of matter \cite{Kitaev2006}. In the gapped phase, the Kitaev model gives rise to itinerant Majorana fermions and non-Abelian Ising anyons in the bulk, accompanied by chiral Majorana modes at the boundary, responsible for phenomena such as the quantized thermal Hall effect \cite{Kitaev2006}. Remarkably, the simplicity and exact solvability of the Kitaev model offer a unique platform for exploring the complex behaviors emerging in quantum spin liquids.

Beyond theoretical interest, the potential realization of the Kitaev model in various two-dimensional materials has spurred considerable experimental efforts%\cite{Jackeli2009, trebst2017kitaev, Winter_2016, Winter2017rev, takagi2019}
. Notably, materials such as $\alpha$-RuCl$_3$ have emerged as prominent candidates for hosting Kitaev physics. However, experimental observations, particularly regarding the interpretation of phenomena like the thermal Hall effect, remain subject to debate, necessitating alternative probes to discern the true nature of the observed states %\cite{Plumb2014,Sandilands2015,Banerjee2016,Baek2017,Kasahara2018,Motome2020,Yokoi_2021,Czajka_2022,Bruin2022}.

In this context, exploring the electrical response of Kitaev spin liquids emerges as a promising avenue for both detection and manipulation of these exotic phases. Theoretical proposals encompass a range of techniques, from electric dipole contributions to subgap optical conductivity to scanning tunneling spectroscopy and interferometry in electrical conductance. Notably, recent work has highlighted the potential for electric polarization and orbital currents associated with localized excitations as signatures of Kitaev spin liquids %\cite{bolens2018,katsura2018,Feldmeier2020,Konig2020,Carrega2020,Udagawa2021,Bauer2023,Aasen2020,halasz2023gatecontrolled,Pereira2020,Banerjee2023}.
 
 

A quantum spin liquid is a fascinating state of matter for magnetic materials \cite{Savary2016}. 
The main property of this state is the lack of spontaneous alignment of the spins, and magnetic order, even at zero temperature. This and other exotic behaviors make them a very interesting topic of research in condensed matter physics. 
The most prominent application for this state of matter are the topological quantum computers \cite{Kitaev2003,Nayak2008}.

The intrinsic property of a quantum spin liquid is a massive degree of entanglement.
To identify this state of matter one should identify (1) topological order, (2) long-range entanglement, and (3) fractionalization. Which means that 
(1) there is a ground state degeneracy that cannot be distinguishable by local measurements, which makes it resistant to errors, 
(2) the global entanglement implies a non-zero topological entanglement entropy \cite{Kitaev_2006,Jiang_2012},
and (3) excitations in the quantum spin liquids carry fractional quantum numbers, and in two-dimensional systems, they can be non-Abelian anyons.



\section{From RVB to Toric code}

... describe the trial wave function described by Anderson in 1973 and how it related to the parton construction (with the Gutzwiller projection)  raised in the context of cuprates and high-Tc superconductivity in 1987

... explain where the topology enters, draw a connection to the Toric code and its ground states

%- in insulating magnets → facilitated by frustrated interaction — corresponding to classical spin systems which display a large ground state degeneracy because the local energy cannot be minimized in a unique way


\section{Kitaev honeycomb model}

The Kitaev model is an exactly solvable spin-1/2 model in a honeycomb lattice which exhibits a quantum spin liquid as a ground state\cite{Kitaev_2006}.  
%The is a very interesting scenario since this phase of matter is very complicated, and in this model one can explore a plenty of properties. 
In the gapped phase, the model exhibits itinerant Majorana fermions and non-Abelian Ising anyons in the bulk and chiral Majorana modes at the boundary. 
The chiral edge modes are responsible for a quantized thermal Hall effect \cite{Kitaev2006}.
The Kitaev model is unique as these complicated properties that emerges in quantum spin liquid phases can be understood in a exactly solvable. 


%The absence of order parameters 
Beyond the natural theoretical interest in the Kitaev model to better understand a quantum spin liquid phase, the model can be approximately realized in different two-dimensional platforms \cite{Jackeli2009, trebst2017kitaev, Winter_2016, Winter_2017, takagi2019}. 
The Jackeli-Khaliullin mechanism works for $t_{2g}$ ions in edge shared octahedra %in which the hopping process assisted by ligand 
%this geometry makes in 
favors a bond-dependent non-diagonal hopping amplitude $t_2$. 
This mechanism happens due to the $90^\circ$ angle in the metal-ligand-metal process, which suppress the isotropic hopping between the $t_{2g}$ three-orbitals of the transition metal 
while favoring an orientating-dependent hopping between different orbitals.
$\alpha\text{-RuCl}_3$ is the currently most prominent candidate Kitaev material \cite{Sandilands2015,Banerjee2016}.
In these materials, sub-dominant interactions beyond the Kitaev model are relevant, which are described by the so-called extended Kitaev model. 
There is some evidence for the observation of the thermal Hall effect in some samples of $\alpha\text{-RuCl}_3$  for a narrow temperature interval \cite{Kasahara2018,Motome2020,Yokoi_2021,Ye2018}, but it is controversial whether this effect is a signature of the  Majorana modes in a chiral quantum spin liquids. For example one could see similar effects for topological magnons \cite{Sahasrabudhe_2020,Czajka_2022} or gapless acoustic phonons \cite{Vinkler_Aviv_2018}.



... few sentences on the controversial experimental realization of Kitaev materials 

\section{PhD Journey}

In view of the difficulty in identifying the QSL phase in Kitaev material, this project arises from the work in Ref.\cite{Pereira2020} which, as well as in other recent theoretical works \cite{Aasen2020,Zhang2020}, proposes an alternative technique using an electrical field to help unambiguously identify this phase. 
 


.. How I started working on this project

what other projects inspired this
the problem we want to solve 
what I have tried 
what is new here


\section{Organization of the thesis}














